\chapter{Agent Evaluation}
\label{ch:evaluation}

\epigraph{``What gets measured gets managed. But what gets measured also gets gamed.''}{}

\learninggoal{
	\item Distinguish between component-level, task-level, and system-level evaluation.
	\item Design evaluation protocols that resist overfitting and contamination.
	\item Understand the major agent benchmarks and their failure modes.
	\item Implement trajectory evaluation beyond simple success/failure metrics.
}

\section{Why Agent Evaluation Is Hard}

Evaluating agents is harder than evaluating \llm s because:
\begin{itemize}
	\item \textbf{Non-determinism:} the same input may produce different trajectories due to sampling and environment interactions.
	\item \textbf{Environment dependence:} success depends on the environment state, which may change between evaluations.
	\item \textbf{Long horizons:} a 20-step agent trajectory cannot be evaluated by checking only the final output---intermediate steps matter.
	\item \textbf{Cost:} each evaluation run costs $10$--$100\times$ more than a single \llm\ prompt.
\end{itemize}

\begin{definition}[Evaluation Levels]
	\begin{itemize}
		\item \textbf{Component evaluation:} test individual modules (tool selection accuracy, memory retrieval precision).
		\item \textbf{Task evaluation:} test end-to-end task completion (success rate, cost, latency).
		\item \textbf{System evaluation:} test the deployed system under realistic load (throughput, reliability, user satisfaction).
	\end{itemize}
\end{definition}

\section{Metrics}

\subsection{Task-Level Metrics}

\begin{longtable}{p{4cm}p{4cm}p{4cm}}
	\toprule
	\textbf{Metric}  & \textbf{Definition}                      & \textbf{Limitation}            \\
	\midrule
	Success rate     & Fraction of tasks completed successfully & Binary, loses partial progress \\
	Goal completion  & Fraction of sub-goals achieved           & Requires sub-goal annotation   \\
	Cost             & Total \llm\ calls + tool calls           & Ignores quality                \\
	Latency          & Wall-clock time to completion            & Hardware-dependent             \\
	Token efficiency & Tokens used per successful task          & Biased toward terse agents     \\
	\bottomrule
\end{longtable}

\subsection{Trajectory-Level Metrics}

Beyond binary success, the quality of the trajectory matters:

\begin{itemize}
	\item \textbf{Optimality ratio:} $\frac{\text{optimal steps}}{\text{agent steps}}$. How many unnecessary steps did the agent take?
	\item \textbf{Backtracking rate:} fraction of steps that are later reverted. High backtracking indicates poor planning.
	\item \textbf{Tool call accuracy:} fraction of tool calls with appropriate arguments.
	\item \textbf{Error recovery rate:} after a tool error, fraction of cases where the agent recovers successfully.
	\item \textbf{Human intervention rate:} how often does the system require human assistance?
\end{itemize}

\begin{equation}
	\text{Optimality ratio} = \frac{|\text{optimal trajectory}|}{|\text{actual trajectory}|}
\end{equation}

\subsection{Efficiency Metrics}

\begin{align}
	\text{Throughput}    & = \frac{\text{completed tasks}}{\text{time}}                                                          \\
	\text{Cost per task} & = \frac{\sum(\text{LLM tokens} \cdot \text{price}) + \sum(\text{tool costs})}{\text{completed tasks}} \\
	\text{Reliability}   & = \frac{\text{tasks without human intervention}}{\text{total tasks}}
\end{align}

\section{Benchmarks}

\subsection{AgentBench}

AgentBench~\cite{liu2023agentbench} evaluates agents across 8 environments:

\begin{longtable}{p{3cm}p{4cm}p{4cm}p{3cm}}
	\toprule
	\textbf{Environment} & \textbf{Task}                   & \textbf{Evaluation} & \textbf{Open?} \\
	\midrule
	Operating system     & Execute shell commands          & Task completion     & Yes            \\
	Web browsing         & Navigate websites, extract info & Goal completion     & Yes            \\
	Database             & Write and run SQL queries       & Query correctness   & Yes            \\
	Knowledge graph      & SPARQL queries                  & Answer accuracy     & Yes            \\
	Digital card game    & Play a strategy game            & Win rate            & Yes            \\
	Household            & Simulated home tasks            & Task completion     & Yes            \\
	\bottomrule
\end{longtable}

\subsection{WebArena}

WebArena~\cite{zhou2023webarena} provides a realistic web environment:

\begin{itemize}
	\item Fully functional web applications (shopping, CMS, forum, git).
	\item 812 tasks requiring navigation, form filling, and content extraction.
	\item Evaluation via: DOM state comparison (did the page change correctly?).
\end{itemize}

Limitation: task diversity is constrained by the pre-built applications.

\subsection{SWE-bench}

SWE-bench~\cite{xiao2024sweb} evaluates code-generation agents on real GitHub issues:

\begin{itemize}
	\item 2,294 task instances from 12 Python repositories.
	\item Task: given an issue description and codebase, produce a patch.
	\item Evaluation: does the patch pass the repository's test suite?
	\item Human performance: $\sim 84\%$ resolved. Best agent (2024): $\sim 50\%$.
\end{itemize}

SWE-bench revealed that many ``successful'' agents were overfit: they exploited patterns in the test set rather than genuinely solving issues.

\section{Evaluation Pitfalls}

\subsection{Contamination}

Agent benchmarks are vulnerable to contamination because:
\begin{itemize}
	\item The \llm\ may have been trained on the benchmark data.
	\item Few-shot examples in the prompt may leak solution patterns.
	\item The agent may search the web and find the answer to a benchmark question.
\end{itemize}

Mitigation: held-out evaluation sets, dynamic tasks (generate new instances procedurally), and output-based evaluation (does the answer work? not does it match the expected answer?).

\subsection{Overfitting to Evaluators}

When using \llm-as-judge for evaluation, the agent may learn patterns that fool the judge:

\begin{itemize}
	\item Verbose answers score higher regardless of correctness.
	\item Self-promotion (``I am very confident about this answer'') inflates scores.
	\item Format manipulation (structuring output to match the judge's expected format).
\end{itemize}

Mitigation: use multiple independent judges, test against held-out rubrics, and periodically validate judge outputs against human annotations.

\subsection{Evaluation Noise}

Agent evaluation is noisy:

\begin{equation}
	\text{Measurement noise} = \sigma_{\text{sampling}} + \sigma_{\text{environment}} + \sigma_{\text{judge}}
\end{equation}

\begin{itemize}
	\item \textbf{Sampling noise:} the agent's policy is stochastic. Run 3--5 trials per configuration.
	\item \textbf{Environment noise:} network delays, API availability. Control for environmental factors.
	\item \textbf{Judge noise:} \llm\ judges are inconsistent. Average over multiple judge calls.
\end{itemize}

\section{Practical Evaluation Design}

\begin{lstlisting}[caption=Agent evaluation pipeline]
  class AgentEval:
      def __init__(self, tasks, n_trials=3):
          self.tasks = tasks
          self.n_trials = n_trials

      def evaluate(self, agent):
          results = []
          for task in self.tasks:
              task_results = []
              for _ in range(self.n_trials):
                  trajectory = agent.run(task.prompt)
                  metrics = self.score(trajectory, task)
                  task_results.append(metrics)
              # Aggregate across trials
              results.append(self.aggregate(task_results))
          return {
              "success_rate": mean(r.success for r in results),
              "avg_cost": mean(r.cost for r in results),
              "avg_latency": mean(r.latency for r in results),
              "per_task": results
          }
\end{lstlisting}

\section{Exercises}

\begin{exercise}
	Design an evaluation protocol for a customer support agent. Define: task taxonomy, success criteria, minimum trials, and evaluation budget.
\end{exercise}

\begin{exercise}
	Implement a trajectory quality scorer for a ReAct agent. Score each trajectory on: correctness, efficiency (steps), backtracking, and error recovery. Validate against human judgments on 50 trajectories.
\end{exercise}

\begin{exercise}
	Identify three sources of evaluation noise in the SWE-bench setup. Propose a protocol to reduce measurement variance.
\end{exercise}

\begin{exercise}
	An agent achieves 90\% on AgentBench but 30\% on a held-out set of similar tasks. What is likely happening? Design a diagnostic experiment.
\end{exercise}
